\section{问题描述}

我们考虑一个在 $\RR^n$ 中由 $N$ 个智能体组成的群体。
令 $\NN=\{1,\dots,N\}$ 表示智能体索引集合。
每个智能体的位置记为 $p_i\in\RR^n$，其动力学由 $m_i$ 阶系统建模
\begin{equation}
    p_i^{(m_i)}=u_i+d_i,
\end{equation}
目标的位置记为 $p_0\in\RR^n$，并记
$p_{ij}=p_i-p_j$，$p_{i0}=p_i-p_0$。

\begin{definition}[围捕性质]
    （改编自 \cite{chenCooperativeTargetfencingProtocol2019}）
    自组织围捕问题被求解，如果
    \begin{enumerate}
        \item[\textbf{P1}] \textbf{凸包围捕：}  $p_0(\infty)\in \mathrm{co}(p(\infty))$；
        \item[\textbf{P2}] \textbf{避碰：}  对给定的安全距离 $d>0$，
              $\| p_{ij}(t) \| > d$，$i, j \in \NN$，$j \neq i$，$\forall ~ t \ge 0$；
        \item[\textbf{P3}] \textbf{速度匹配（刚性编队）：}
              $\lim_{t\to\infty}\|v_i-v_0\|=0$，$\forall i\in\NN$。
    \end{enumerate}

    这里，
        $v_i=\dot p_i$，$v_0=\dot p_0$，且
        $
        \mathrm{co}(p) = \left\{ \sum_{i=1}^N \lambda_i p_i
        : \lambda_i \geq 0, \forall ~ i \in \NN ~ \textrm{且} \sum_{i=1}^N \lambda_i = 1 	\right\}
        $
        是由所有智能体张成的凸包，
        $\mu(S)$ 是集合 $S$ 的勒贝格测度。
        在 $\RR^2$ 中它是凸多边形的面积，在 $\RR^3$ 中它是多面体的体积。
\end{definition}



\subsection{APF 函数}

下文尝试中使用的排斥型 APF 项为
\[
    \phi_i^r = \sum_{j \in \NN, j\neq i} \alpha_r(\| p_{ij} \|)
                                     \frac{p_{ij}}{\|p_{ij}\|}
    = \sum_{j \in \NN_i(t)} \alpha_r(\| p_{ij} \|)
                                     \frac{p_{ij}}{\|p_{ij}\|},
\]
其中
\[
    \NN_i(t)=\{j\in\NN\setminus\{i\}:\|p_{ij}(t)\|<\mu\}
\]
是激活邻域集合。
连续函数 $\alpha_r: (d, \infty) \to [0, \infty)$ 满足
\begin{enumerate}
    \item[A1.] $\lim_{s \to d^+} \int_s^\mu \alpha_r(\sigma)\mathrm{d}\sigma = +\infty$；
    \item[A2.] $\alpha_r(s) = 0$，$\forall ~ s \in [\mu, \infty)$。
\end{enumerate}
下文中我们令 $\phi_i=\phi_i^r$。
由于成对项是反对称的，在安全集合
$\{p:\|p_{ij}\|>d,\ i\ne j\}$ 上有 $\sum_{i=1}^N\phi_i=0$。

\subsection{APF 函数的微分}

为说明为何直接对原始 APF 项求导是不可取的，考虑一个邻域集合固定且
所有成对距离均大于 $d$ 的时间区间。
在本段中，额外假设 $\alpha_r$ 在 $(d,\mu)$ 上是 $C^2$ 的；
这比 A1--A2 更强，仅用于展示形式导数。
令
\[
    v_{ij}:=\dot p_{ij},\quad
    a_{ij}:=\ddot p_{ij},
\]
\[
    r_{ij}:=\|p_{ij}\|,
\]
并定义
\[
    \beta(r):=\frac{\alpha_r(r)}{r}.
\]
则
\[
    \phi_i
    =\sum_{j\in\mathbb N_i}\alpha_r(r_{ij})
      \frac{p_{ij}}{r_{ij}}
    =\sum_{j\in\mathbb N_i}\beta(r_{ij})p_{ij}.
\]
对每一对定义雅可比矩阵
\[
    A_{ij}
    :=\beta(r_{ij})I
      +\frac{\beta'(r_{ij})}{r_{ij}}p_{ij}p_{ij}^{T}.
\]
则一阶导数可简洁地写作
\[
    \dot\phi_i
    =\sum_{j\in\mathbb N_i}A_{ij}v_{ij}.
\]
二阶导数已明显更复杂：
\[
    \ddot\phi_i
    =\sum_{j\in\mathbb N_i}
    \left(A_{ij}a_{ij}+\dot A_{ij}v_{ij}\right),
\]
其中，记 $r=r_{ij}$，$p=p_{ij}$，$v=v_{ij}$ 且
$\dot r=p^T v/r$，
\[
\begin{aligned}
    \dot A_{ij}
    =&\ \beta'(r)\dot r\, I
      +\frac{\beta'(r)}{r}
       \left(vp^T+pv^T\right)\\
     &+\frac{\beta''(r)r-\beta'(r)}{r^2}
       \dot r\,pp^T .
\end{aligned}
\]
因此，计算 $\dot\phi_i$ 需要所有激活邻域的相对速度以及导数 $\alpha_r'$。
计算 $\ddot\phi_i$ 进一步需要相对加速度、$\alpha_r''$ 以及
相对位置与速度的乘积。
对于高阶反步，仍需要 $\phi_i$ 的更高阶导数。

这在控制器中难以使用，原因如下。
首先，邻域集合 $\mathbb N_i(t)$ 在成对距离跨越影响半径时会发生变化，
因此 $\phi_i$ 通常仅分段光滑。
其次，在安全边界 $r_{ij}=d$ 附近，$\alpha_r$ 及其导数可能变得奇异，
使得 $\dot\phi_i$ 和 $\ddot\phi_i$ 在数值上病态或无界。
第三，$\ddot\phi_i$ 依赖于相对加速度，其中包含其他智能体的控制和干扰；
这并非纯局部的位置信息。
最后，要求这些导数与 APF 围捕的原始动机（即基于主要相对位置的
简单分布式律）相冲突。